% =====================================================================
% Chapter 5, part 2: Conclusions. Each paragraph carries the objectives
% of Section 1 (sec:objectives) it closes, in parentheses.
% =====================================================================

\section{Conclusions}
\label{sec:conclusions}

This thesis set out to determine whether Parkinson's disease can be detected
from the micro-Doppler gait signatures of a validated clinical radar
network, and to do so under an evaluation a clinician could trust: every
subject scored by a model that had never seen them. Five conclusions follow
from the work; each notes, in parentheses, the objectives of
Section~\ref{sec:objectives} it closes.

\textbf{A complete, pre-registered pipeline was built and frozen before the
results were known} (O2, O3). Preprocessing, windowing, model families,
hyperparameters, validation and statistics were fixed in advance
(Chapter~\ref{ch:methods}), so every number reported in Chapter~\ref{ch:results} is the first and only
answer the protocol gave, not the best of many attempts.

\textbf{Under subject-independent validation, no classifier was
distinguishable from recording duration alone} (O1, O5). Classical models, a compact
CNN, the ImageNet probe, a sequence model on velocity envelopes, and
the field's own AlexNet all land in a band whose confidence intervals
contain the duration floor of 0.610 (Sections~\ref{sec:res-baselines},
\ref{sec:res-deep} and~\ref{sec:disc-alexnet}). The attention analysis and the
background-suppression experiment attribute the remaining score to
per-recording character rather than to gait shape
(Sections~\ref{sec:disc-attention} and~\ref{sec:disc-intervention}).

\textbf{The gap to the literature was reproduced, not just alleged} (O5).
With the validation design relaxed to the field's random hold-out, the same data
and networks produce subject-level AUCs up to 0.89 that collapse on people
excluded from training (Sections~\ref{sec:disc-protocol}
and~\ref{sec:disc-alexnet}). The published
numbers answer a recognition question; the clinical question demands
subject-independent evaluation, and this thesis provides, to the best of our
knowledge, the first such evaluation for radar-based Parkinson's
classification.

\textbf{One clinically sensible signal survives} (O1). Patients take measurably
longer, with the excess concentrated in standing up and turning; radar
measures this timing reliably, which amounts to an automated
timed-up-and-go measurement (Section~\ref{sec:disc-positive}). It is a
modest signal, but it is real, interpretable, and it is exactly what the
source system was validated to measure.

\textbf{The negative result is the contribution} (O4). This work demonstrates,
on validated clinical hardware, precisely how radar gait classifiers come to
look better than they are, provides the protocol that prevents it, and
states which signals in this dataset are exhausted and which were never
testable. Whatever signal richer data holds, the standard defended here will
measure it honestly, on people the model has never seen.
